% SPDX-License-Identifier: CC-BY-4.0
\documentclass[11pt,letterpaper]{article}

\usepackage[T1]{fontenc}
\usepackage{amsmath,amssymb}
\usepackage{tikz}
\usepackage{pgfplots}
\pgfplotsset{compat=1.18}
\usepackage{algorithm}
\usepackage{algpseudocode}
\usepackage{booktabs}
\usepackage{enumitem}
\usepackage[hidelinks]{hyperref}

\newcommand{\N}{\mathbb{N}}
\newcommand{\Z}{\mathbb{Z}}
\newcommand{\Q}{\mathbb{Q}}
\newcommand{\R}{\mathbb{R}}
\newcommand{\C}{\mathbb{C}}
\newcommand{\Map}[3]{#1 \colon #2 \to #3}
\newcommand{\Mod}[1]{#1\text{-}\mathrm{Mod}}
\newcommand{\GroupHom}{\mathrm{Hom}_{\mathbf{Grp}}}
\newcommand{\RingHom}{\mathrm{Hom}_{\mathbf{Ring}}}

\newtheorem{theorem}{Theorem}
\newtheorem{lemma}[theorem]{Lemma}
\newtheorem{definition}[theorem]{Definition}
\newtheorem{remark}[theorem]{Remark}
\newenvironment{proof}{\noindent\textit{Proof.}}{\hfill\(\square\)}

\title{Generalized Collatz Library}
\author{ascendnoosphere}
\date{\today}

\begin{document}
\maketitle

\section*{Positive additive all-seed outcomes}

Throughout, fix integers $c\ge1$ and $D\ge2$. The multiplier is exactly $q=1$.
On the positive integers define the original raw map
\[
 C(n)=\begin{cases}n/D,&D\mid n,\\ n+c,&D\nmid n.\end{cases}
\]
Division has precedence. Write $C^k$ for $k$ original instructions, and put
\[
 g=\gcd(c,D),\quad L=D/g,\quad B=c(L-1),\quad
 N=\left\lfloor\frac{B}{D-1}\right\rfloor.
\]
The \href{../../QCD.Additive.001/tex/Additive.tex}{access and core theorems}
will be used with their exact raw clocks. Their needed statements are
recalled below. The
\href{../../QCD.Additive.002/tex/AdditiveCycles.tex}{cycle classification}
identifies all positive periodic states by whole modular orbits; it does
not identify basin membership from the initial residue alone.

\section{Legal events and stopped traps}

\begin{lemma}[First division and finite core]\label{LemmaAccess}
If $g\mid n$, there is a unique integer $\ell(n)$ with
$0\le\ell(n)<L$ and $D\mid n+\ell(n)c$. Every preceding affine instruction is
legal, and
\[
 H(n)=\frac{n+\ell(n)c}{D}=C^{\ell(n)+1}(n)>0.
\]
If $g\nmid n$, then no division is possible and $C^j(n)=n+jc$ for every
$j\ge0$. For accessible $n>N$, one has $H(n)<n$. For accessible
$1\le n\le N$, one has $1\le H(n)\le N$. In particular, every positive seed
reaches a positive core state or an inaccessible state.
\end{lemma}

\begin{proof}
Write $c=gc_0$ and $D=gL$, where $\gcd(c_0,L)=1$. If $g\mid n$, the
congruence $\ell c_0\equiv-n/g\pmod L$ has exactly one solution in
$\{0,\ldots,L-1\}$. No smaller nonnegative index can give a division,
by uniqueness. The quotient is positive because its positive numerator
is divisible by $D$. If $g\nmid n$, divisibility of $n+jc$ by $D$ would
force $g\mid n$, a contradiction. Finally
\[
 DH(n)=n+\ell(n)c\le n+B.
\]
For $n>N$ the inequality $(D-1)n>B$ implies $H(n)<n$.
For $n\le N$, one has $n+B<D(N+1)$, since $B<(D-1)(N+1)$;
the integral quotient is therefore at most $N$. Iterating while above the
core and accessible strictly decreases a positive integer. This iteration
terminates; it does not require a bound on an unknown infinite orbit.
\end{proof}

For the finite-state argument define a stopped event model
\[
 e(n)=\begin{cases}H(n),&g\mid n,\\ n,&g\nmid n,\end{cases}
 \qquad
 w(n)=\begin{cases}\ell(n)+1,&g\mid n,\\0,&g\nmid n.\end{cases}
\]
The stopped value is not a raw fixed point. Its zero clock records that the
event path has stopped at an affine trap. Define accumulated raw clocks by
\[
 K_0=0,\qquad K_{i+1}=K_i+w(e^i(n)).
\]

\begin{lemma}[Exact event recovery]\label{LemmaRecovery}
For every $i\ge0$, one has $e^i(n)=C^{K_i}(n)$ and $0\le K_i\le iL$.
The clocks are nondecreasing. An accessible phase contributes a strictly
positive increment. A positive core state remains in $\{1,\ldots,N\}$ under
$e$, including after stopping.
\end{lemma}

\begin{proof}
For an accessible phase, Lemma~\ref{LemmaAccess} gives
$e(x)=C^{w(x)}(x)$ and $1\le w(x)\le L$. For an inaccessible phase,
$e(x)=x=C^0(x)$ and $w(x)=0$. Induct on $i$, using
$C^{a+b}(n)=C^b(C^a(n))$, to recover the same endpoint at the same clock.
Summing $0\le w\le L$ gives the upper bound and monotonicity.
Core invariance follows from Lemma~\ref{LemmaAccess} in the accessible case
and from $e(x)=x$ in the stopped case.
\end{proof}

\section{Finite core outcomes and recovery}

\begin{theorem}[Core return or true trap]\label{TheoremCore}
For every $1\le r\le N$, there is $a\le NL$ such that either
\[
 g\nmid C^a(r),
\]
or there is a positive integer $p\le NL$ with
\[
 C^{a+p}(r)=C^a(r).
\]
\end{theorem}

\begin{proof}
The $N+1$ event states $r,e(r),\ldots,e^N(r)$ all belong to the $N$-element
positive core. If some state $e^i(r)$, $i\le N$, is inaccessible, choose
$a=K_i\le NL$ and use Lemma~\ref{LemmaRecovery}.
Otherwise the pigeonhole principle gives $0\le i<j\le N$ with
$e^i(r)=e^j(r)$. In this second case the phase $e^i(r)$ is accessible, so
$K_{i+1}>K_i$, and clock monotonicity yields $K_j>K_i$.
Choose the common witness
\[
 a=K_i,\qquad p=K_j-K_i.
\]
It satisfies $0<p\le NL$ and $a\le NL$; recovery gives the displayed raw
return. A repeated stopped state is never used as a positive-period witness.
If $N=0$ there is no positive $r\le N$, so no base case has been omitted.
\end{proof}

\begin{theorem}[All-seed dichotomy]\label{TheoremOutcome}
For every positive seed $n$, exactly one of the following alternatives holds:
\begin{enumerate}
\item There are integers $u\ge0$ and $p\ge1$ such that
$C^{u+p}(n)=C^u(n)$; the trajectory is eventually periodic.
\item There is $u\ge0$ such that $g\nmid C^u(n)$ and
\[
 C^{u+j}(n)=C^u(n)+jc\quad(j\ge0).
\]
The entire raw sequence tends to infinity.
\end{enumerate}
\end{theorem}

\begin{proof}
Lemma~\ref{LemmaAccess} supplies one entry witness $v$, with
$r=C^v(n)>0$ either in the core or inaccessible. In the latter case the
exact affine tail begins at $u=v$. In the core case apply
Theorem~\ref{TheoremCore} to this same $r$. If it produces a trap at clock
$a$, put $u=v+a$. If it produces a return of length $p$ at clock $a$, use
the same $u=v+a$ and the iteration identity to obtain
$C^{u+p}(n)=C^u(n)$. Thus witnesses from separate trajectories are never
combined.

The affine tail tends to infinity because $c>0$; removing its finite prefix
does not change the full-sequence limit. Conversely a returned phase $r$
recurs at all times $u+kp$, $k\ge0$. These times tend to infinity and their
values stay equal to $r$, excluding a full-sequence limit of infinity.
Determinism makes the finite raw block of length $p$ repeat forever.
The alternatives are therefore exhaustive and disjoint.
\end{proof}

\begin{theorem}[Uniform explicit witness bound]\label{TheoremBudget}
Define
\[
 \delta(x)=(D-1)x-B,\qquad \epsilon=\delta(N+1).
\]
Then $1\le\epsilon\le D-1$. For every integer $h\ge0$ with
$\delta(n)<D^h\epsilon$, the witness $u$ in Theorem~\ref{TheoremOutcome}
can be chosen with
\[
 u\le(h+N)L,
\]
and a return witness can be chosen with $0<p\le NL$.
The nonoptimal choice $h=n$ always satisfies the hypothesis.
\end{theorem}

\begin{proof}
Write $B=(D-1)N+r$ with $0\le r<D-1$. Then
$\epsilon=D-1-r$, proving its bounds. For every accessible event,
\[
 D\delta(H(x))\le\delta(x),
\]
as follows by substituting $DH(x)\le x+B$.
If the first $h$ events all remain accessible and above the core, iteration
gives $D^h\delta(H^h(n))\le\delta(n)$. But an above-core endpoint has
$\delta(H^h(n))\ge\epsilon$, contradicting the strict budget hypothesis.
Thus the core or a trap is reached after at most $h$ events and at raw time
$v\le hL$. If entry is into the core, Theorem~\ref{TheoremCore} supplies
$a\le NL$ and the same-witness recovery uses $u=v+a$.
This proves the bounds without estimating an unknown transient separately.
For $h=0$, the strict defect condition already places a positive seed in the
core; no event needs to be taken.

Finally Bernoulli's inequality gives
\[
 \delta(n)\le n(D-1)<1+n(D-1)\le D^n\le D^n\epsilon.
\]
Hence $h=n$ is an explicit finite choice. This is a bound in the numeric
seed and parameters, not a uniform polynomial bit-complexity assertion.
\end{proof}

\section{Exact basin and parameter criteria}

\begin{theorem}[Periodicity criteria]\label{TheoremCoprime}
A positive seed is eventually periodic if and only if
$g\mid C^t(n)$ for every $t\ge0$. Consequently,
\[
 \text{every positive seed is eventually periodic}
 \quad\Longleftrightarrow\quad \gcd(c,D)=1.
\]
For a fixed positive seed, full-sequence escape occurs if and only if a
finite inaccessible-state witness is reached.
\end{theorem}

\begin{proof}
Suppose a raw return occurs at clocks $u,u+p$, with $p>0$, and a trap is
reached at any clock $t$. Applying $C^t$ to the return identity gives
$C^{u+p+t}(n)=C^{u+t}(n)$. Both times are after the trap, so its exact affine
tail gives two values differing by $pc>0$, a contradiction. Thus a periodic
tail cannot coexist with a trap at any clock. The converse follows from the
exhaustive dichotomy: if every raw state is accessible, the trap alternative
is excluded. If $g=1$, every raw state of every positive seed is accessible.
If all positive seeds are eventually periodic, apply the criterion to seed
$1$ at clock zero to obtain $g\mid1$, hence $g=1$.
The escape equivalence follows from the exact affine tail and the fact that
a returned phase cannot escape, as proved in Theorem~\ref{TheoremOutcome}.
\end{proof}

\begin{theorem}[No-cycle criterion]\label{TheoremNoCycle}
Every positive seed escapes if and only if there is no positive raw periodic
state. Therefore the complete finite cycle test of
\href{../../QCD.Additive.002/README.md}{QCD.Additive.002} decides universal
escape when its whole-orbit enumeration is complete.
\end{theorem}

\begin{proof}
A positive periodic state cannot escape. Conversely, if a seed does not
escape, Theorem~\ref{TheoremOutcome} produces a positive returned phase,
which is a positive raw periodic state. The companion cycle theorem
exhausts such states by whole orbits in its explicit finite modular window.
A truncated enumeration or a count from an incomplete search is not used
in this recovery argument. Finite orbit and order methods are discussed
in \cite{DummitFoote04}.
\end{proof}

\subsection*{Separating examples and boundaries}
For $(c,D)=(9,6)$, one has $N=1$ but the companion window is empty;
therefore every positive seed escapes despite the nonempty core.
For $(c,D)=(10,4)$, the supported decomposition is $c=2\cdot5$ and the
window radius is $1$. Its only possible representative is $1$, but its
successor under multiplication by $4^{-1}\equiv4\pmod5$ is $4$.
No whole orbit fits, so again every positive seed escapes.
For $(c,D)=(6,4)$, the raw cycle is $2\mapsto8\mapsto2$, whereas the initially
accessible seed $4$ divides to $1$ and then follows $1+6j$.
Initial access is not a sufficient basin criterion.

All escape claims require $c>0$. At $c=0$, inaccessible states are raw fixed
points, not escaping tails; the homogeneous unit-multiplier theorem applies.
Several event and finite-budget lemmas allow this zero-offset arithmetic
boundary, but the positive-offset dichotomy and its escape criteria do not.
No assertion here concerns a multiplier other than $q=1$.

\section{Finite certificates and bounded computation}

The Go implementation stores arbitrary-precision seed, phase, raw prefix
clock, and positive return clock. An event consumes $\ell+1$ original
instructions, even when this number is too large for machine integers.
A trap certificate records its inaccessible phase and zero return interval.
A periodic certificate records one prefix and a positive closed event
interval, not a certificate of least period.
The verifier recomputes legal first-division events without using the
classifier's finite-core history. It checks the same prefix phase and clock,
and then checks either true inaccessibility or the exact positive raw return.
For example, $c=1$, $D=2^{200}$, and seed $1$ yield one event whose raw return
period is $2^{200}$; verification uses its congruence, not literal replay of
all affine instructions.

The classifier has separate event and core-storage limits. Their failure
produces no new outcome certificate. A zero-event inaccessible state can
still be certified with a zero event budget. Dense core storage is needed
only after accessible core entry; every concurrent worker owns its cache
and arithmetic workspace. Numeric integer fields are distinct from
machine-sized event counts and resource limits. The bounded raw-reference
tests are computational checks, not a replacement for the all-seed proof.

\bibliographystyle{amsalpha}
\bibliography{AdditiveOutcome}
\end{document}
